\section{Plaquette model details}

The minimal plaquette model used in the main text is
\begin{equation}
H =
J \sum_{\langle i,j\rangle} Z_i Z_j
+
h \sum_i X_i
+
\lambda Z_0 Z_1 Z_2 Z_3.
\end{equation}
This model is intended to isolate the structural ingredients needed for finite EDT dynamics: local couplings, transverse control that enables operator spreading, and a loop term that introduces genuinely nonlocal correlations. It is not presented as a quantitatively fitted Hamiltonian for K$_3$C$_{60}$.

The accompanying reproducible implementation is the package \texttt{src/plaquette\_model}. The Hamiltonian constructor is \texttt{build\_hamiltonian(J, h, lambda\_loop, n\_qubits=4)}. Exact state trajectories are generated by \texttt{exact\_time\_evolution}, OTOCs by \texttt{compute\_otoc}, nearest-neighbor ZZ correlators by \texttt{compute\_two\_point\_correlations}, subsystem von Neumann entropies by \texttt{compute\_entanglement\_entropy}, and the finite-regime composite score by \texttt{compute\_edt\_dynamical\_score}. The first-order circuit-inspired approximation is exposed by \texttt{trotter\_evolution(J, h, lambda\_loop, dt, n\_steps)}.

The dynamical score is designed to be finite-regime rather than monotonic. It combines normalized OTOC growth, average nearest-neighbor correlation strength, and an entropy factor that is maximal at intermediate normalized entropy and suppressed both near zero and near saturation. This reflects the intended interpretation of the minimal model: the relevant EDT regime is the interval where spreading, correlation retention, and structured entanglement are jointly effective.
